\section{Recursive procedures and unconditional termination}
\label{sec:algorithm}

There are four mutually recursive procedures. The main procedure
$\Solve$ narrows a box, using repeated smaller-rank solves through
$\Ask$. It chooses the direction of the final narrowing through
$\OrientPlus$, which repeats $\Orient$. We give all instructions
first and then prove that this recursion terminates on every tape.
The later correctness argument will place probabilistic guarantees on
these already defined computations.

The graph, transformed operator $F$, and number $B_*$ are fixed for the
whole run. All vector operations are coordinatewise. Scalar shifts such
as $a+C$ use the same scalar at every coordinate, and rational floors,
including floors of negative values, are taken toward minus infinity.
Set
\begin{equation}\label{eq:global-parameters}
\begin{gathered}
 \eta=\frac45,\qquad
 g=1+\left\lceil\log_2(B_*+1)\right\rceil,\qquad
 s=\min\{2^j:j\in\N,\ 2^j\ge2^{17}(g+1)\},\qquad h=\frac1s,\\
 J=2[s(n+1)]^6+1,\qquad \varepsilon=2^{-J}.
\end{gathered}
\end{equation}
The repetition count $J$ is odd. Only the analysis uses $\varepsilon$;
the algorithm never constructs its denominator $2^J$.

A permitted solver input is
\[
 (a,b;B,\alpha,\beta,k,m),\qquad
 a,b\in\Z^V,\quad a\le b,\quad1\le B\le B_*,\quad b-a\le B,
\]
with positive rational support weights and budgets. Unmentioned
parameters are inherited by a subcall. A changed parameter belongs
only to that subcall and never alters its caller's parameters.
Each subcall uses fresh fair bits, independent of its pre-call history.
Repetitions on a fixed input are independent.

The deterministic routine $\Basic(a,b)$ is exactly that of
Lemma~\ref{lem:box-solution}: start at $p=a$, compute
$p'=\clip_{[a,b]}F(p)$ synchronously, return $p$ when $p'=p$, and
otherwise replace $p$ by $p'$. It will be invoked only at small
reference widths.

\subsection{The two repetition wrappers}

The following wrappers refer to procedures specified immediately below.
Their termination, together with the rest of the mutual recursion, is
proved in Proposition~\ref{prop:termination}.

\begin{procedure}[Repeated comparison]\label{proc:ask}
For its given solver input, $\Ask$ runs $J$ independent copies of
$\Solve$ and returns their coordinatewise median.
\end{procedure}

\begin{procedure}[Repeated orientation]\label{proc:orient-plus}
For its given orientation input, $\OrientPlus$ runs $J$ independent
copies of $\Orient$. It returns the majority choice in $\{Y,Z\}$
and the concatenation of all exceptional-set lists in trial order.
It retains minority-trial lists and duplicate entries.
\end{procedure}

Every $\Ask$ below reduces a budget or applies
Lemma~\ref{lem:reweighting}. The main continuation of $\Solve$ is
a single unamplified call. Orientation is used only at $R\ge1$, and
its own width recursion leaves all support parameters unchanged.

\subsection{Narrowing the box}

The tightening loop raises the lower bound or lowers the upper bound
whenever a repeated solve moves some coordinate by at least $\delta$.
After it stops, the choice $Z$ caps the upper boundary, while $Y$ caps
the lower boundary. The exceptional-set calls supplement that one
main continuation.

\begin{procedure}[The solver]\label{proc:solve}
Given $(a,b;B,\alpha,\beta,k,m)$, compute $R,\mu$ from
Equation~\eqref{eq:rank-parameter} and execute these steps.
\begin{enumerate}[label=\textbf{S\arabic*.},ref=S\arabic*,leftmargin=*]
\item\label{alg:solve-base}
If $R<1$, return the vector with coordinate $b_i$ when $\beta_i>m$
and $a_i$ otherwise. If $R\ge1$ but $B\le128s$, return $\Basic(a,b)$.

\item\label{alg:solve-tighten}
Define $\delta=\lfloor B/(4s)\rfloor$ and $C=B-\delta$.
Run the following loop with fixed $B,\alpha,\beta,k,m$ and changing
bounds, still denoted by $a,b$.
\begin{enumerate}[label=(\alph*),leftmargin=*]
\item Evaluate
\[
 r=\Ask(a,b;B,\alpha,\beta,\eta k,m).
\]
If $r_i-a_i\ge\delta$ for some $i$, replace the entire lower-bound
vector $a$ by $r$ and restart the loop. Otherwise proceed without
changing $a$.
\item Evaluate
\[
 r=\Ask(a,b;B,\alpha,\beta,k,\eta m).
\]
If $b_i-r_i\ge\delta$ for some $i$, replace the entire upper-bound
vector $b$ by $r$ and restart. Otherwise stop the loop with both
bounds unchanged in this pass.
\end{enumerate}

\item\label{alg:solve-orient}
Using the final tightening box, evaluate
\[
 (\chi,\mathcal H)=\OrientPlus(a,b;B,\alpha,\beta,k,m),
\]
where $\chi\in\{Y,Z\}$ and $\mathcal H$ is a list of subsets of $V$.

\item\label{alg:solve-z}
When $\chi=Z$:
\begin{enumerate}[label=(\alph*),leftmargin=*]
\item For every entry $H$ of $\mathcal H$, form
\[
 \alpha_i^H=\begin{cases}\alpha_i/2,&i\in H,\\\alpha_i,&i\notin H,\end{cases}
 \qquad r_H=\Ask(a,b;B,\alpha^H,\beta,3k/5,m).
\]
\item Execute the main continuation
\[
 p=\Solve\bigl(a,\min(b,a+C);C,\alpha,\beta,k,m\bigr),
\]
and return the coordinatewise maximum of $p$ and all $r_H$.
\end{enumerate}

\item\label{alg:solve-y}
When $\chi=Y$:
\begin{enumerate}[label=(\alph*),leftmargin=*]
\item For every entry $H$ of $\mathcal H$, form
\[
 \beta_i^H=\begin{cases}\beta_i/2,&i\in H,\\\beta_i,&i\notin H,\end{cases}
 \qquad r_H=\Ask(a,b;B,\alpha,\beta^H,k,3m/5).
\]
\item Execute the main continuation
\[
 p=\Solve\bigl(\max(a,b-C),b;C,\alpha,\beta,k,m\bigr),
\]
and return the coordinatewise minimum of $p$ and all $r_H$.
\end{enumerate}
\end{enumerate}
\end{procedure}

A listed call always uses the final tightening box, independent of all
previous listed outputs. Duplicate entries cause distinct calls. The
maximum or minimum always includes $p$, so an empty list is permitted.

\subsection{Choosing a direction from moving centers}

An orientation trial repeatedly averages results from the two reduced
budgets. Large final displacement can determine its choice immediately.
Otherwise, a fresh sampled iterate supplies the center of a smaller
orientation call.

\begin{procedure}[Orientation]\label{proc:orient}
For $(a,b;B,\alpha,\beta,k,m)$ with $R\ge1$, compute $\mu$ and proceed
as follows.
\begin{enumerate}[label=\textbf{O\arabic*.},ref=O\arabic*,leftmargin=*]
\item\label{alg:orient-base}
If $B\le128s$, compute $p=\Basic(a,b)$ and put
\[
 P=\{i:p_i\ge a_i+B/2\}.
\]
Return $Y$ with the one-entry list $(V\setminus P)$ if $\mu(P)>1/2$;
otherwise return $Z$ with $(P)$.

\item\label{alg:orient-loop}
Initialize
\[
 q_0=\left\lfloor\frac{a+b}{2}\right\rfloor,\qquad
 D=\left\lfloor\frac{B}{32s}\right\rfloor,\qquad
 E_0=\left\lfloor\frac B8\right\rfloor.
\]
For $j=0,\ldots,s-1$, make independent calls
\[
\begin{aligned}
 r_1&=\Ask(q_j-D,q_j+D;2D,\alpha,\beta,\eta k,m),\\
 r_2&=\Ask(q_j-D,q_j+D;2D,\alpha,\beta,k,\eta m),
\end{aligned}
\]
and set
\[
 q_{j+1}=\clip_{[a,b]}\left(\left\lfloor\frac{r_1+r_2}{2}\right\rfloor\right).
\]

\item\label{alg:orient-motion}
Form the disjoint displacement sets
\[
 P=\{i:(q_s)_i-(q_0)_i>2hB\},\qquad
 N=\{i:(q_s)_i-(q_0)_i<-2hB\}.
\]
If $\mu(P)>1/2$, return $Y$ and $(V\setminus P)$.
If that test fails but $\mu(N)>1/2$, return $Z$ and $(V\setminus N)$.

\item\label{alg:orient-child}
If neither test returned, draw $j$ uniformly from $\{0,\ldots,s-1\}$
using $\log_2s$ fresh bits. Evaluate
\[
 (\chi,\mathcal H)=
 \Orient(q_j-E_0,q_j+E_0;2E_0,\alpha,\beta,k,m).
\]
Return $\chi$ and the list $\mathcal H$ followed by $P,N$.
\end{enumerate}
\end{procedure}

Store the iterates $q_0,\ldots,q_{s-1}$ as they are produced so that the
sampled one is available. This is finite local storage and will be charged
in Section~\ref{sec:complexity}. The comparison cores used to analyze
these procedures are never computed by the algorithm.

\subsection{Output bounds and termination}
\label{alg:termination-subsection}

Before discussing accuracy, we prove that every execution is finite and
that all recursive inputs and outputs have their stated types. These
facts must survive failed comparisons: a failed tape may change the
answers, but it must not invalidate the program or its resource bounds.
Recall that $d(R)$ is the least nonnegative integer satisfying
$(9/10)^{d(R)}R<1$.

\begin{proposition}\label{prop:termination}
For every permitted $\Solve$ input, the execution and all its subroutine
executions terminate on every random tape. The outputs of $\Solve$
and $\Ask$ are integer vectors in their respective input boxes.
The outputs of $\Orient$ and $\OrientPlus$ have the stated choice
and list types, and each listed set has input measure at most $1/2$.
Moreover:
\begin{enumerate}[label=(\roman*)]
\item Tightening makes at most $8ns$ updates and $16ns+2$ calls to
$\Ask$.
\item A chain of main $\Solve$ continuations contains at most
\begin{equation}\label{eq:main-chain-length}
 K=8s(g+1)+1
\end{equation}
solver invocations.
\item One $\Orient$ trial visits at most $g+1$ orientation invocations
and returns at most $2g+2$ sets. Consequently $\OrientPlus$ returns
at most $J(2g+2)$ sets, counting repetitions.
\end{enumerate}
\end{proposition}

\begin{proof}
We induct on the solver rank and establish the output assertions at the
same time as termination. At rank zero, $R<1$, and the endpoint rule
returns immediately in the given box. Fix a positive rank and assume
these properties for every smaller-rank solver input.

An ordinary budget change multiplies $R$ by $4/5$. Reweighting a set
with $\mu(H)\le1/2$ multiplies it by
$(3/5)(1+\mu(H))\le9/10$. For either change,
\[
 (9/10)^{d(R)-1}R_{\mathrm{child}}
 \le(9/10)^{d(R)}R<1,
\]
so the child has smaller rank. Such $\Ask$ calls terminate by the
induction hypothesis. Their medians remain integer vectors in their
input boxes, irrespective of comparison accuracy.

\paragraph{Orientation at the current rank.}
The only solver calls within orientation are through budget-reduced
$\Ask$. At width $B\le128s$, $\Basic$ terminates by
Lemma~\ref{lem:box-solution}; the returned set is either $P$ of
measure at most $1/2$ or the complement of a set of measure greater
than $1/2$. At larger width, both
$D=\floor{B/(32s)}$ and $E_0=\floor{B/8}$ are positive. Each reduced
call uses a legitimate centered box with reference width
$2D\le B/(16s)$. Its result lies between $q_j-D$ and $q_j+D$.
The integer floor of their average lies between these same integer
endpoints. Since $q_j\in[a,b]$, subsequent clipping keeps the next
center in $[a,b]$ and cannot increase its distance from $q_j$. Therefore
\begin{equation}\label{alg:center-step-bound}
 |(q_{j+1})_i-(q_j)_i|\le D\qquad(i\in V).
\end{equation}
This proves that all $s$ loop steps terminate with valid centers.

The early returns again use complements of sets with measure greater
than $1/2$. When both tests fail, the appended $P,N$ each have measure
at most $1/2$. The child box has positive width $2E_0\le B/4$ and
unchanged support parameters, hence unchanged $\mu$. A separate
induction on this decreasing positive integer width proves orientation
termination and the measure assertion. After $r$ child edges its width
is at most $B_*/4^r$. Since $g>\log_2 B_*$, there are at most $g$
edges, or $g+1$ invocations. The final one returns one set and each
ancestor adds two, giving $1+2g\le2g+2$ entries. Repeating this finite
trial $J$ times gives a well-defined majority because $J$ is odd, and
concatenation proves the list bound for $\OrientPlus$.

\paragraph{Tightening at the current rank.}
The solver's small-width branch is already covered by basic iteration.
Otherwise $B>128s$, and
\begin{equation}\label{alg:delta-bound}
 \delta=\floor{B/(4s)}\ge\frac{B}{8s}>0,
 \qquad0<C=B-\delta<B.
\end{equation}
Every tightening output is inside the current box. A lower update only
raises $a$, an upper update only lowers $b$, and both preserve $a\le b$.
At least one coordinate moves by $\delta$. The integer quantity
$\sum_i(b_i-a_i)$ starts at most $nB$ and decreases by at least
$\delta$ at each update. There are at most $nB/\delta\le8ns$ updates.
A pass contains at most two $\Ask$ calls and either restarts after an
update or is the final pass. This gives $2(8ns+1)=16ns+2$ calls.

\paragraph{The main continuation and its returned vector.}
Orientation now terminates by the argument just given. All of its sets
have the required measure on every tape, so every reweighted call has
smaller rank and terminates. The sole main successor has a valid subbox
and reference width $C$. Along a chain of such successors,
\[
 B_{r+1}\le B_r\left(1-\frac1{8s}\right),\qquad
 B_r\le B_*\exp\left(-\frac r{8s}\right).
\]
After $8s(g+1)$ edges this last bound would be below one, because
$g+1>\ln B_*$. In fact a chain stops already at width at most $128s$.
Thus it has at most $K$ invocations. This reasoning uses the reference
width, so it remains valid when a cap leaves the actual box unchanged.

Work backwards along this finite chain. All reweighted outputs lie in
the node's final tightening box, and the main output lies in its
specified subbox. Their coordinatewise maximum or minimum is still in
the tightening box, itself contained in the node's input box. Backward
induction therefore proves the solver output assertion and termination
at the current rank. No correctness event has been used at any stage.
Since orientation is invoked only after the test $R<1$ fails, all uses
of positive rank above apply to the actual calls. This completes the
rank induction.
\end{proof}

